% -------------------------- HEAD --------------------------
\documentclass[a4paper,11pt]{article}
\usepackage[T1]{fontenc}				% codifica dei font
\usepackage[utf8]{inputenc}				% lettere accentate da tastiera
\usepackage[english]{babel}				% lingua del documento
\usepackage{lipsum}						% per generare testo fittizio (placeholder text)

% Load the custom style file
\usepackage{AndreaStyle}
% The file `AndreaStyle.sty` is stored in: `D:\Programmi e Applicazioni\texlive\texmf-local\tex\latex\local` for Windows.
% The file `AndreaStyle.sty` is stored in: `/usr/local/texlive/texmf-local/tex/latex/local` for Ubuntu (desktop).
% This won't work in Overleaf, until the AndreaStyle.sty file is added to the project

\begin{document}
	\title{Ringraziamenti}
	\maketitle
	
	Introduzione 
	
	\textcolor{red}{Introduzione: chi mi conosce lo sa. Tendo a scrivere molto. Scrivere questa tesi è stato un lavoro lungo e difficile, ma non è niente rispetto alla difficoltà nello scrivere queste righe.}
	
	{\color{red} Date le premesse iniziali e alcune fasi dello svolgimento del lavoro, non era affatto scontato che, alla fine, sarei stato soddisfatto della tesi che ho scritto. 
	Invece posso affermarlo, e il merito va all'impeccabile assistenza che ho ricevuto. Questa lettera è per ringraziare tutti quelli che vi hanno contribuito.}
	
	\textcolor{red}{Vorrei scrivere qualcosa di denso e significativo, che rispecchi quanto profondamente gli anni di studio hanno inciso su di me ma anche, e soprattutto, quanto siano sentiti i ringraziamenti. Nel fare ciò, mi confronto con il dilemma: da una parte vorrei scrivere qualcosa di quanto più ``distaccato'' possibile -- nel tentativo di rendere questi ringraziamenti ``eterni'', dall'altra mi sento che instanziarli al contesto presente è l'unico modo per renderli davvero ``memorabili''.}
	\bigskip  
	
	La prima persona da ringraziare è senza dubbio il mio relatore, Federico Giovanni Poloni. 
	
	La ringrazio per tutto il supporto che mi ha dato, non solo durante lo svolgimento della tesi, ma anche prima, per i miei ultimi esami. La ringrazio per la sua pazienza e per la meticolosità con cui ha corretto le mie lunghissime bozze. La ringrazio per avermi sempre risposto prontamente ed esaustivamente.
	
	Mi rendo conto che tutto ciò non è scontato, purtroppo. Non potevo veramente chiedere di più da un relatore, confesso però che avevo capito ben presto che lei si attiene a standard molto alti. 
	
	Il poter contare su tutto questo è stato molto importante per il mio lavoro di tesi. Se questa è stata un'esperienza proficua da cui ho imparato molto -- e non solo per ciò che riguarda l'esponenziale di matrice e l'aritmetica di macchina -- lo si deve soprattutto alla sua professionalità e alle sue doti umane. 
	
	Grazie, soprattutto, per avermi motivato e per aver creduto in me, anche per tutto ciò che riguarda il post-tesi.
	\medskip  
	
	Dedica a nonno 
	
	{\color{red}Vorrei dedicare questa tesi a mio nonno Cosimo. Mi ricordo sempre, nonno, quando negli ultimi giorni ti chiamavo e ti raccontavo dei bei voti che avevo preso a scuola. Sentivo che era altro ciò di cui dovevamo parlare; nonostante volessi chiederti altro, nonostante sapessi che non ne avrei avuto l'occasione per molto tempo, non riuscivo. Tutto ciò che riuscivo a fare era allontanare il pensiero di non averti più parlando di studio. Purtroppo, non ho smesso.
	
	Mi ricordo che quando ti dicevo che avevo preso un bel voto rispondevi dicendo che dovevo sempre andare avanti così. Beh, eccomi qui. Ho dato del mio meglio. Spero che tu sia orgoglioso di me.}
	\medskip 
	
	Grazie famiglia
	
	{\color{red}Un ringraziamento speciale va alla mia famiglia.
	Grazie a mamma. Grazie papà. Grazie Adri. Grazie Ale.}
	\medskip 
	
	Grazie famiglia tutta
	
	{\color{red}Ci tengo a ringraziare anche la mia famiglia allargata. Siete i migliori. Non potevo sperare di avere una famiglia migliore.
	
	In particolare, grazie a nonna Rosaria. Sei veramente d'esempio per me. Un riferimento molto importante.}
	\medskip 
	
	Grazie Tribù
	
	{\color{red} Grazie alla Tribù. Spesso mi figuro scenari passati alternativi, a come sarebbe la mia vita se avessi fatto una scelta piuttosto che un'altra. Non riesco a pensare a un universo parallelo in cui non ho incontrato voi. Siete la cosa migliore che mi sia mai capitata.
		
	Un grazie particolare ai membri della Tribù che conosco dalle elementari: Lorenzo, Akira, Matteo. Ogni volta che penso a questa cosa non riesco a credere che questo miracolo sia successo davvero.
	
	Grazie a membri singoli della tribù: Marco, Gian, Pietro, Art, Ark, Luci, Nene \textbf{due parole per tutti dai}.}
	\medskip 
	
	Grazie amici singoli
	
	{\color{red} Un grazie al Dime, a Renato, (rimetto Marco?), insomma agli ``amici singoli''.}
	\medskip 
	
	Grazie Sfascio
	
	{\color{red} Un grazie agli amici del gruppo ``Sfascio''. Mi avete dato quella tranquillità, quella stabilità di cui avevo bisogno. La felicità sta anche nel prendersi una pizza tutti insieme.}
	\medskip 
	
	Grazie gente AL
	
	{\color{red} Grazie agli amici dell' AL Lucca, o Maghi delle Mura, o Kraken Lucca, o Roku Lucca, sicuramente sto scordando qualche nome\dots, insomma grazie a tutti gli amici conosciuti tramite i giochi di ruolo. Con voi ho condiviso esperienze importanti, ma soprattutto ho condiviso una passione che mi ha accompagnato per lungo tempo e di cui farò tesoro per sempre.
	
	In particolare grazie a Marco, Roberta, Leo, Rob, Enri, \textbf{Altri dai. Altri? O nessuno?}}
	\medskip 
	
	Grazie NetRiders
	
	{\color{red} Ci tengo molto a ringraziare i NetRiders + Torquatizers: Giordano, Andrea, Yuri, DavideB, Angelo, DavideM. \`E stato un onore oltre che un piacere collaborare con voi e condividere la vostra amicizia. Sono contento che si sia creata una forte identità di gruppo. Spero che vi abbia lasciato anche solo un briciolo di quel che ha lasciato a me.
		
	Grazie per aver reso questi anni universitari così divertenti. Siete il gruppo di compagni di corso perfetto, siete ciò che rimpiango di non aver avuto prima, nella triennale. 
	
	Grazie per avermi sopportato. Mi mancherà il non preparare più esami con voi.}
	\medskip
	
	Altri ringraziamenti qui sicuramente\dots
	\bigskip
	
	Ringraziamenti finali
	
	{\color{red} Infine, Grazie a Eleonora per avermi costretto a imparare come aprire, grazie al Guido per avermi costretto a imparare come chiudere, e grazie a Martina per avermi costretto a imparare tutto ciò che sta nel mezzo.}
	
\end{document}